\chapter{Audio output \textsuperscript{Neumayer}}
\label{chap:Audio output}

One of the most important aspects of any audio-related project is how well the software 
performs real-time audio output operations. Low latency is especially critical 
when users expect an instant response of the application. 

\section{Differences between platforms}

While Qt allows multi-platform development of natively compiled applications, 
there are still cases where differences between supported platforms have to be 
taken into account. It was decided that a timer-based approach was inappropriate 
and therefore, a thread-based solution would have to be used. The problem with basing audio output on a periodic timer is that, while it worked perfectly on Linux, Qt on Windows 
didn't provide timers with sufficiently high accuracy. This is due to a deficiency in the Windows audio stack which Qt's timers are based on.

\section{Audio output using Qt}

The Qt framework provides facilities to access audio hardware in an abstract way. 
In the case of Beat The Score the used classes are {\tt QAudioFormat}, {\tt QAudioDeviceInfo} and {\tt QAudioOutput}.

\subsection{QAudioFormat}

{\tt QAudioFormat} is used to tell the audio hardware information about the sound data that will be 
written to it.
\\

\begin{lstlisting}
QAudioFormat format;
format.setSampleSize(16);
format.setSampleRate(44100);
format.setChannelCount(1);
format.setCodec("audio/pcm");
format.setByteOrder(QAudioFormat::LittleEndian);
format.setSampleType(QAudioFormat::SignedInt);
\end{lstlisting}

Audio data is made of individual samples. The sample size defines how many bits a single 
sample has. The sample rate describes how many samples audio data can have within the 
time span of one second. Channel count refers to the number of concurrent channels the 
audio data has. In Beat The Score only audio data with a single channel is used. 
\\
The codec tells the audio device which format the to-be-expected audio data has. 
According to the Qt documentation \footnote{\label{foot:1}http://qt-project.org/doc/qt-5/qaudiodeviceinfo.html\#supportedCodecs}, 
'All platform and plugin implementations should provide support for: "audio/pcm"'.
\\
Beat The Score uses audio files in the RIFF WAVE format, which is an uncompressed, 
raw audio format with meta data about its sample rate, bits per sample, et cetera. 
According to the RIFF WAVE specification \footnote{\label{foot:2}https://ccrma.stanford.edu/courses/422/projects/WaveFormat/} samples are expected to be both little-endian and, in the case of 16-bit samples, signed.

\subsection{QAudioDeviceInfo}

In order to access the audio device with a supported audio format Qt provides 
functionality to check if the audio hardware actually supports the declared format.

\begin{lstlisting}
QAudioDeviceInfo info(QAudioDeviceInfo::defaultOutputDevice());
if (!info.isFormatSupported(format)) {
    cout << "format sadly unsupported, trying to find nearest format" << endl;
    format = info.nearestFormat(format);
}
\end{lstlisting}

The static public method {\tt QAudioDeviceInfo::defaultOutputDevice()} is a factory method that, 
depending on the targeted platform, returns an object which represents the audio hardware 
to output sound. The source code of Qt that is referred to in this section is released 
as version 5.2.1 of Qt. \footnote{\label{foot:3}https://qt.gitorious.org/qt/qtmultimedia/source/\\9a16423610405c0329fe40235313212946f08b05:src/multimedia/audio}
On systems where the ALSA audio architecture is used (Linux and Android), 
the first audio device that can be found on the system is used, 
according to the source code of qaudiodeviceinfo\_alsa\_p.cpp. 
On Windows (qaudiodeviceinfo\_win32\_p.cpp) Qt selects the default 
audio device based on the WAVE\_MAPPER device identifier.
\\
The output device is provided to the constructor of {\tt QAudioDeviceInfo} which allows 
information about the device to be gathered. This is used to decide whether the 
device actually supports the desired audio format. In case the requested format is 
not supported, the application tries to find a format that is similar to the requested 
format using the {\tt nearestFormat()} method of {\tt QAudioDeviceInfo}. This method compares 
the formats that are supported by the device with the provided format and returns the 
most similar format, based on codec, byte order, sample type, sample size, channel count, 
and sample rate.

\subsection{QAudioOutput}

After an audio device with a supported format has been selected it is possible to create an 
object of type {\tt QAudioOutput} to send sample data to the audio hardware. 
According to the Qt documentation \footnote{\label{foot:5}http://qt-project.org/doc/qt-5/audiooverview.html}, {\tt QAudioOutput} provides two modes to accomplish this: 

\begin{description}
\item [Pull mode] The {\tt start()} method of a {\tt QAudioOutput} object can be given a pointer to {\tt QIODevice} object as a parameter. In this case the {\tt QIODevice} contains a buffer with mixed audio samples. Doing this will make the {\tt QAudioOutput} decide for itself when to pull audio samples from the {\tt QIODevice}. The decision when to pull samples is based on how many samples can be sent to the sound card for one period. The duration of a period can vary during runtime. \hfill \\
\item [Push mode] When the {\tt start()} method of a {\tt QAudioOutput} object is called without a parameter, it returns a pointer to a {\tt QIODevice} object. The application has to write audio samples to this {\tt QIODevice} by itself periodically. This approach is more appropriate for real-time audio as the fluctuating nature of the pull mode does not guarantee low latency sound output. \hfill \\
\end{description}

It was decided to use the push mode in Beat The Score as it is more appropriate for real-time audio use cases.

\section{Pitching audio samples}

To easily and quickly pitch audio samples the application uses the SoundTouch library by 
Olli Parviainen \footnote{\label{foot:5}http://www.surina.net/soundtouch}. SoundTouch provides an implementation 
of the Synchronous-Overlap-Add algorithm. 
\\ \\
To quote an article of Olli Parviainen \cite{soundtouch1}: 
\linebreak
\begin{quote}
SOLA, being short-hand for Synchronous-OverLap-Add, works by cutting the sound data into 
shortish sequences between tens to hundreds milliseconds each, and then joining these 
sequences back together with a suitable amount of sound data either skipped or partially 
repeated in between these sequences, to achieve a shorter or longer playback time than originally.
\end{quote}

For details on how to use the SoundTouch API please refer to section \ref{sec:SoundTouch} for more information.
\\
Beat The Score ships with samples for guitars, pianos and drums in the RIFF WAVE format. 
The file name of these samples corresponds to its MIDI note number as defined 
by the MIDI specification. 

\begin{center}
\begin{tabular}{ | l | l | }
  \hline
  MIDI note number & note \\ \hline
  .. & .. \\ \hline
  48 & C \\ \hline
  49 & C\# \\ \hline
  50 & D \\ \hline
  51 & D\# \\ \hline
  52 & E \\ \hline
  \hline
\end{tabular}
\end{center}

For each of the instruments there are only sample files of the note C in different octaves. 
When the application starts, it reads these existing samples from the file system and 
pitches missing semitone steps to get valid samples for each octave in a loop until 
there is a sample for every possible MIDI note number.
\\ \\
Internally, the application creates a vector of {\tt NoteSample}s where a 
{\tt NoteSample} is a C++ data structure.
\\ \\ \\ \\ \\ \\ \\ \\ 
\begin{lstlisting}
struct NoteSample {
    NoteSample()
    {
        buffer = NULL;
    }
    ~NoteSample()
    {
        if (buffer != NULL)
            delete buffer;
    }

    char * buffer;
    int bufferSize;
};

vector<shared_ptr<NoteSample>> notes; // index is the note number as specified in the MIDI spec
\end{lstlisting}

A {\tt NoteSample} represents exactly one note. What MIDI note number a {\tt NoteSample} 
relates to is identified by the vectors index. The {\tt char *} inside the 
struct {\tt NoteSample} points to an address in memory where its audio samples are 
stored and the buffer size defines the buffers size in memory. As a result 
there are 128 {\tt NoteSample}s in the vector for the MIDI note numbers from 0 to 127.

\section{Audio mixing}

An audio mixer is needed to play multiple notes at the same time. 
This is implemented in the C++ class called {\tt AudioBuffer}. 
The header file of {\tt AudioBuffer} looks as follows:
\\
\begin{lstlisting}
class AudioBuffer
{

private:
    vector<shared_ptr<NoteSample>> noteBuffers;
    vector<int> notePositions;
    QMutex mutex;

public:
    AudioBuffer(qint64 maxSize);
    ~AudioBuffer();
    qint64 readData(char * mixedBuffer, qint64 size);
    void writeData(shared_ptr<NoteSample> note);
    void writeData(shared_ptr<NoteSample> note, int offset);
    void clear();
};
\end{lstlisting}

The vector {\tt noteBuffers} contains shared pointers to {\tt NoteSample}s and \\ 
{\tt notePosition} contains the information on how many samples of the {\tt NoteSample}s 
buffer have been read. 
The {\tt QMutex} object is provided by the Qt framework and used to handle access to the {\tt NoteSample}s, so that {\tt NoteSample}s cannot be added to the mixer while audio is being mixed at the same time.
\\
The method {\tt readData} is used to mix the buffers of each {\tt NoteSample} together. \\
\begin{lstlisting}
qint64 AudioBuffer::readData(char * mixedBuffer, qint64 size)
{
    QMutexLocker(&this->mutex);
    // Let's only add up samples that overlap at a single point in time.
    // Also, do mixing in relation to the longest buffer in the vector.
    int longestSample = 0;

    for(int i = 0; i < noteBuffers.size(); i++) {
        if(noteBuffers[i]->bufferSize > longestSample) {
            longestSample = noteBuffers[i]->bufferSize < size ?
                        noteBuffers[i]->bufferSize : size;
        }
    }

    // Step through every single point in "time" and mix samples that overlap
    for(int i = 0; i < longestSample; i+=2) {
        short mixedValue = 0;
        int mixBufferCount = 0; // How many buffers can be mixed at this point in time?
        for(int j = 0; j < noteBuffers.size(); j++) {
            shared_ptr<NoteSample> sample = noteBuffers[j];
            if(sample != NULL && notePositions[j] >= 0 &&
                    (sample->bufferSize - notePositions[j]) > 0) {
                mixBufferCount++;
            }
        }

        for(int j = 0; j < noteBuffers.size() && mixBufferCount > 0; j++) {
            shared_ptr<NoteSample> sample = noteBuffers[j];
            if(sample != NULL && notePositions[j] > 0) {

                // Can this part of the sample be mixed?
                int length = (sample->bufferSize - notePositions[j]) > size ?
                        size : (sample->bufferSize - notePositions[j]);

                if(length > 0) {
                    short tmpVal = 0;
                    memcpy(&tmpVal, &(sample->buffer[notePositions[j]]), 2);
                    tmpVal = tmpVal / mixBufferCount;
                    mixedValue += tmpVal;
                }
            }

            notePositions[j] += 2;
        }

        memcpy(&mixedBuffer[i], &mixedValue, 2);
    }

    // Remove notes from the vector if they have been mixed completely
    for(int i = noteBuffers.size()-1; i >= 0; i--) {
        shared_ptr<NoteSample> sample = noteBuffers[i];
        if(sample->bufferSize < notePositions[i]) {
            noteBuffers.erase(noteBuffers.begin()+i);
            notePositions.erase(notePositions.begin()+i);
        }
    }

    // Append silence in case there is still some room for values
    if(longestSample < size) {
        memset(mixedBuffer+longestSample, 0, size-longestSample);
    }

    return size;
}
\end{lstlisting}

The parameter {\tt size} describes how many "parts" of each individual buffer have to be mixed. 
These parts are generally referred to as samples. After the desired samples have been read 
from a buffer the corresponding position is incremented by the value of the {\tt size} parameter. 
When mixing, the numerical values of concurrent samples are added together and divided by 
the number of concurrent {\tt NoteSample}s. The resulting value gets copied into the provided 
character pointer called {\tt mixedBuffer}. If the position of a {\tt NoteSample} exceeds its 
buffer size the {\tt NoteSample} and its position are removed from their respective vectors. 
After the process of mixing is done the method returns the provided size.
\\
The method {\tt writeData} is used to put {\tt NoteSample}s into the vector and 
initializing its position, either by initializing the {\tt position} with 0 
or by an optionally provided offset. The optional {\tt offset} paramenter describes by how many samples 
mixing of this specific {\tt NoteSample} has to be delayed. A position of 0 indicates that the 
{\tt NoteSample}s buffer has to be mixed as soon as the {\tt readData()} method is called, 
a position of negative value delays the mixing of the specific {\tt NoteSample} 
into the future and a position of positive value would cause {\tt readData()} to skip samples.

\section{Output of instruments}

Catching MIDI signals and turning them into something that the application can 
understand is done by the {\tt MidiInput} class. Depending on the platform the 
application runs on, this class uses either RtMidi on Windows and Linux or libusb on 
Android to read the MIDI data and putting its values in a {\tt Note} object. 

\section{Output of background music}

Background music is considered to be every track of the music except the selected input 
track. These tracks are looped through in a separate background thread inside the {\tt NoteOutput}  
object. Looping through the tracks is done periodically, just like the actual audio 
output itself, every 20 milliseconds. This background thread runs synchronously with the 
game's internal clock to not mix samples into the master buffer when the game pauses or 
when the user is in the main menu. 